% GENERATED FILE. Edit canonical lessons, then run npm run generate:books.
% canonical-source-hash: fnv1a64:c182dd0813648c2e
% canonical-lessons: KA-C157-velapatti, KA-C157-suchane, KA-C157-badige, KA-C157-jahiratu, KA-C157-rasidi, KA-R157-first-pass-saying-no-and-looking-back, KA-R157-second-pass-looking-ahead-and-reading-signs

\chapter{A Timetable, a Notice, the Rent, an Advert, a Receipt}
\label{ch:ka-a-timetable-a-notice-the-rent-an-advert-a-receipt}
\hlchaptermodality{157}

\begin{chapteropening}
\textbf{By the end of this chapter:} \emph{I can read a timetable, a notice, an advert and a receipt, and act on what they say.}

Complete the last lesson of chapter 157: I can read a timetable, a notice, an advert and a receipt, and act on what they say.
\end{chapteropening}

\section[\texorpdfstring{\kn{ವೇಳಾಪಟ್ಟಿ} (vēḷāpaṭṭi)}{ವೇಳಾಪಟ್ಟಿ (vēḷāpaṭṭi)}]{\kn{ವೇಳಾಪಟ್ಟಿ} (vēḷāpaṭṭi) — a timetable}

\label{lesson:KA-C157-velapatti}



\begin{quote}
Before the new one: say the Kannada for to decide, then the Kannada for to expect.
\end{quote}

\subsection*{You'll want to know: \kn{ವೇಳಾಪಟ್ಟಿ}}
\textbf{\kn{ವೇಳಾಪಟ್ಟಿ}} — \emph{vēḷāpaṭṭi} — \textquotedblleft{}a timetable\textquotedblright{}. Read the board: \textbf{\kn{ಬಸ್}: \kn{ಬೆಳಿಗ್ಗೆ} \kn{೮}} — \emph{bas: beḷigge eṇṭu} — \textquotedblleft{}Bus: 8 a.m.\textquotedblright{} The bus leaves at eight, so be there before eight.

\begin{grammarlens}[title={What you've built}]
One more word for this chapter.
\end{grammarlens}

\subsection*{Your turn}
\begin{itemize}
\raggedright
  \item \emph{Say it:} \emph{vēḷāpaṭṭi}
  \item \emph{Say it:} \emph{vēḷāpaṭṭi}, once more
  \item \emph{Say it:} say \emph{vēḷāpaṭṭi}
  \item \emph{Recall:} say the Kannada for to decide, then the Kannada for to expect, then say \emph{vēḷāpaṭṭi} again
\end{itemize}

\begin{tcolorbox}[breakable,colback=teal!4,colframe=teal!35!black,title={Before you move on}]
What does \kn{ವೇಳಾಪಟ್ಟಿ} mean? (\textquotedblleft{}A timetable\textquotedblright{}.) Say it once more.
\end{tcolorbox}
\section[\texorpdfstring{\kn{ಸೂಚನೆ} (sūcane)}{ಸೂಚನೆ (sūcane)}]{\kn{ಸೂಚನೆ} (sūcane) — a notice, an instruction}

\label{lesson:KA-C157-suchane}



\begin{quote}
Before the new one: say the Kannada for to expect, then the Kannada for a timetable.
\end{quote}

\subsection*{You'll want to know: \kn{ಸೂಚನೆ}}
\textbf{\kn{ಸೂಚನೆ}} — \emph{sūcane} — \textquotedblleft{}a notice, an instruction\textquotedblright{}. Read the notice: \textbf{\kn{ಇಂದು} \kn{ಅಂಗಡಿ} \kn{ಮುಚ್ಚಿದೆ}} — \emph{indu aṅgaḍi muccide} — \textquotedblleft{}The shop is closed today.\textquotedblright{} Come back on another day.

\begin{grammarlens}[title={What you've built}]
One more word for this chapter.
\end{grammarlens}

\subsection*{Your turn}
\begin{itemize}
\raggedright
  \item \emph{Say it:} \emph{sūcane}
  \item \emph{Say it:} \emph{sūcane}, once more
  \item \emph{Say it:} say \emph{sūcane}
  \item \emph{Recall:} say the Kannada for to expect, then the Kannada for a timetable, then say \emph{sūcane} again
\end{itemize}

\begin{tcolorbox}[breakable,colback=teal!4,colframe=teal!35!black,title={Before you move on}]
What does \kn{ಸೂಚನೆ} mean? (\textquotedblleft{}A notice, an instruction\textquotedblright{}.) Say it once more.
\end{tcolorbox}
\section[\texorpdfstring{\kn{ಬಾಡಿಗೆ} (bāḍige)}{ಬಾಡಿಗೆ (bāḍige)}]{\kn{ಬಾಡಿಗೆ} (bāḍige) — rent; for hire}

\label{lesson:KA-C157-badige}



\begin{quote}
Before the new one: say the Kannada for a timetable, then the Kannada for a notice.
\end{quote}

\subsection*{You'll want to know: \kn{ಬಾಡಿಗೆ}}
\textbf{\kn{ಬಾಡಿಗೆ}} — \emph{bāḍige} — \textquotedblleft{}rent; for hire\textquotedblright{}. The money for a house: \textbf{\kn{ಮನೆ} \kn{ಬಾಡಿಗೆ}} — \emph{mane bāḍige} — \textquotedblleft{}the rent for the house\textquotedblright{}.

\begin{grammarlens}[title={What you've built}]
One more word for this chapter.
\end{grammarlens}

\subsection*{Your turn}
\begin{itemize}
\raggedright
  \item \emph{Say it:} \emph{bāḍige}
  \item \emph{Say it:} \emph{bāḍige}, once more
  \item \emph{Say it:} say \emph{bāḍige}
  \item \emph{Recall:} say the Kannada for a timetable, then the Kannada for a notice, then say \emph{bāḍige} again
\end{itemize}

\begin{tcolorbox}[breakable,colback=teal!4,colframe=teal!35!black,title={Before you move on}]
What does \kn{ಬಾಡಿಗೆ} mean? (\textquotedblleft{}Rent; for hire\textquotedblright{}.) Say it once more.
\end{tcolorbox}
\section[\texorpdfstring{\kn{ಜಾಹೀರಾತು} (jāhīrātu)}{ಜಾಹೀರಾತು (jāhīrātu)}]{\kn{ಜಾಹೀರಾತು} (jāhīrātu) — an advertisement}

\label{lesson:KA-C157-jahiratu}



\begin{quote}
Before the new one: say the Kannada for a notice, then the Kannada for rent.
\end{quote}

\subsection*{You'll want to know: \kn{ಜಾಹೀರಾತು}}
\textbf{\kn{ಜಾಹೀರಾತು}} — \emph{jāhīrātu} — \textquotedblleft{}an advertisement\textquotedblright{}. Read the advert: \textbf{\kn{ಮನೆ} \kn{ಬಾಡಿಗೆಗೆ} \kn{ಇದೆ}} — \emph{mane bāḍigege ide} — \textquotedblleft{}House for rent.\textquotedblright{}

\begin{grammarlens}[title={What you've built}]
One more word for this chapter.
\end{grammarlens}

\subsection*{Your turn}
\begin{itemize}
\raggedright
  \item \emph{Say it:} \emph{jāhīrātu}
  \item \emph{Say it:} \emph{jāhīrātu}, once more
  \item \emph{Say it:} say \emph{jāhīrātu}
  \item \emph{Recall:} say the Kannada for a notice, then the Kannada for rent, then say \emph{jāhīrātu} again
\end{itemize}

\begin{tcolorbox}[breakable,colback=teal!4,colframe=teal!35!black,title={Before you move on}]
What does \kn{ಜಾಹೀರಾತು} mean? (\textquotedblleft{}An advertisement\textquotedblright{}.) Say it once more.
\end{tcolorbox}
\section[\texorpdfstring{\kn{ರಸೀದಿ} (rasīdi)}{ರಸೀದಿ (rasīdi)}]{\kn{ರಸೀದಿ} (rasīdi) — a receipt}

\label{lesson:KA-C157-rasidi}



\begin{quote}
Before the new one: say the Kannada for rent, then the Kannada for an advertisement.
\end{quote}

\subsection*{You'll want to know: \kn{ರಸೀದಿ}}
\textbf{\kn{ರಸೀದಿ}} — \emph{rasīdi} — \textquotedblleft{}a receipt\textquotedblright{}. Ask for it when you pay: \textbf{\kn{ದಯವಿಟ್ಟು} \kn{ರಸೀದಿ} \kn{ಕೊಡಿ}} — \emph{dayaviṭṭu rasīdi koḍi} — \textquotedblleft{}The receipt, please.\textquotedblright{}

\begin{grammarlens}[title={What you've built}]
One more word for this chapter.
\end{grammarlens}

\subsection*{Your turn}
\begin{itemize}
\raggedright
  \item \emph{Say it:} \emph{rasīdi}
  \item \emph{Say it:} \emph{rasīdi}, once more
  \item \emph{Say it:} say \emph{rasīdi}
  \item \emph{Recall:} say the Kannada for rent, then the Kannada for an advertisement, then say \emph{rasīdi} again
\end{itemize}

\begin{tcolorbox}[breakable,colback=teal!4,colframe=teal!35!black,title={Before you move on}]
What does \kn{ರಸೀದಿ} mean? (\textquotedblleft{}A receipt\textquotedblright{}.) Say it once more.
\end{tcolorbox}
\section[(dialogue)]{Saying no and looking back — a first pass}

\label{lesson:KA-R157-first-pass-saying-no-and-looking-back}



\begin{quote}
Say the words for an advertisement and a receipt.
\end{quote}

\subsection*{Your turn}
Say these aloud:

\begin{itemize}
\raggedright
  \item nobody, never, nowhere, anyone, anything
  \item recently, already, last, once, to spend
\end{itemize}

\begin{tcolorbox}[breakable,colback=teal!4,colframe=teal!35!black,title={Before you move on}]
Nobody? (\textbf{\kn{ಯಾರೂ}}, \emph{yārū}.) Never? (\textbf{\kn{ಎಂದಿಗೂ}}, \emph{endigū}.) Nowhere? (\textbf{\kn{ಎಲ್ಲಿಯೂ}}, \emph{elliyū}.) Anyone? (\textbf{\kn{ಯಾರಾದರೂ}}, \emph{yārādarū}.) Anything? (\textbf{\kn{ಏನಾದರೂ}}, \emph{ēnādarū}.) Recently? (\textbf{\kn{ಇತ್ತೀಚೆಗೆ}}, \emph{ittīcege}.) Already? (\textbf{\kn{ಈಗಾಗಲೇ}}, \emph{īgāgalē}.) Last? (\textbf{\kn{ಕಳೆದ}}, \emph{kaḷeda}.) Once? (\textbf{\kn{ಒಮ್ಮೆ}}, \emph{omme}.) To spend? (\textbf{\kn{ಕಳೆ}}, \emph{kaḷe}.)
\end{tcolorbox}
\section[(dialogue)]{Looking ahead and reading signs — a second pass}

\label{lesson:KA-R157-second-pass-looking-ahead-and-reading-signs}



\begin{quote}
Say the words for next and from now on.
\end{quote}

\subsection*{Your turn}
Say these aloud:

\begin{itemize}
\raggedright
  \item next, from now on, the future, to decide, to expect
  \item a timetable, a notice, rent, an advertisement, a receipt
\end{itemize}

\begin{tcolorbox}[breakable,colback=teal!4,colframe=teal!35!black,title={Before you move on}]
Next? (\textbf{\kn{ಮುಂದಿನ}}, \emph{mundina}.) From now on? (\textbf{\kn{ಇನ್ನು} \kn{ಮುಂದೆ}}, \emph{innu munde}.) The future? (\textbf{\kn{ಭವಿಷ್ಯ}}, \emph{bhaviṣya}.) To decide? (\textbf{\kn{ನಿರ್ಧರಿಸು}}, \emph{nirdharisu}.) To expect? (\textbf{\kn{ನಿರೀಕ್ಷಿಸು}}, \emph{nirīkṣisu}.) A timetable? (\textbf{\kn{ವೇಳಾಪಟ್ಟಿ}}, \emph{vēḷāpaṭṭi}.) A notice? (\textbf{\kn{ಸೂಚನೆ}}, \emph{sūcane}.) Rent? (\textbf{\kn{ಬಾಡಿಗೆ}}, \emph{bāḍige}.) An advertisement? (\textbf{\kn{ಜಾಹೀರಾತು}}, \emph{jāhīrātu}.) A receipt? (\textbf{\kn{ರಸೀದಿ}}, \emph{rasīdi}.)
\end{tcolorbox}
